%%==========================================================================
%% CONTRIBUTION 1
%%==========================================================================
\begin{frame}[shrink=5]{C1 Motivation: Curriculum-Based Team Scaling}
\vspace{2em}
\textbf{Motivation:}
\begin{itemize}
    \item Inspired by Smit et al.'s~\footcite{smit2023} 
    work using MARL to scale training in 2v2 to playing 11v11 in a football-like game.
\end{itemize}
\begin{columns}
    \begin{column}{0.5\textwidth}
        \textbf{Method:}
        \begin{enumerate}
            \item Trained teams with 4 methods
            \begin{itemize}
                \item Naive Heuristic
                \item ``Team Work'' Heuristic
                \item 2v2 Trained Agents
                \item 11v11 Trained Agents
            \end{itemize}
            \item Evaluated through full games (11v11) against each other.
        \end{enumerate}
    \end{column}
    \begin{column}{0.43\textwidth}
        \begin{figure}
            \centering
            \includegraphics[width=0.45\linewidth]{smit2v2.png}
            \includegraphics[width=0.45\linewidth]{smit_football.png}
            \caption{Smit et al.'s Game~\cite{smit2023}, 2v2 and 11v11 versions.}
            \label{table:smit_envs}
        \end{figure}
    \end{column}
\end{columns}
\vspace{0.5em}
\textbf{Key Result:}
\begin{itemize}
    \item Computational limitations led to generalization issue and overall underperformance.
\end{itemize}

\note[note]{During Lit review foudn Smit}
\note[note]{Inspired by their austere conditions compare to}
\note[note]{Contemporeneous to bigger successes like OpenAI Five and AlphaStar}
\note[note]{\textbf{Consequence:}}
\note[note]{emphasis on efficiency and ease of access}

\note[note]{\textbf{Their method:}}
\note[note]{four training conditions}
\note[note]{one eval condition}

\note[note]{Unimpressive results, but showed some promise}
\end{frame}

\begin{frame}{C1 Research Questions}
    % \begin{block}{Research Questions}
        \begin{itemize}
            \item[\textbf{RQ1.1}] Can smaller-team pretraining + policy duplication +
                retraining improve training efficiency (agent-steps)
                without sacrificing final performance?
            \item[\textbf{RQ1.2}] How does effectiveness vary across environments with
                different heterogeneity forms?
        \end{itemize}
    % \end{block}

\note[note]{Recall research questions}
\note[note]{Taking up where Smit left off}
\note[note]{but hoping to improve results}
\note[note]{AND extend to heterogeneous systems}
\end{frame}

\begin{frame}{C1 Methodology}
    \begin{columns}
        \begin{column}{0.45\textwidth}
            \begin{block}{Training pipeline:}
                \begin{enumerate}
                    \item \textbf{Phase 1 — Pretrain:} Train a small team ($k$ agents)
                        until a checkpoint criterion is reached.
                    \item \textbf{Phase 2 — Upsample \& Retrain:} Duplicate policies to
                        construct a target team of $N$ agents; continue training
                        with all policies updating independently.
                \end{enumerate}
            \end{block}
        \end{column}
        \begin{column}{0.5\textwidth}
            \begin{figure}
                \centering
                \includegraphics[width=0.95\linewidth]{training_1.tex}
                \caption{Resampling between training phases.}
            \end{figure}
        \end{column}
    \end{columns}

\note[note]{New approach building on theirs}
\note[note]{We add retraining after the upscaling portion}
\note[note]{We examing different ratios of pre-post training}
\note[note]{We do a greater number of team size changes}
\note[note]{Neither of us use parameter sharing}
\end{frame}

\begin{frame}[shrink=5]{The Agent-Steps Metric}
\vspace{3em}
\begin{columns}
    \begin{column}{0.52\textwidth}
        \textbf{Why not other common metrics?}
        \begin{itemize}
            \item \textbf{Training steps} obscure the cost difference
                between 2-agent and 8-agent runs.
            \item \textbf{Wall-clock time} varies with hardware and cluster load.
        \end{itemize}

        \vspace{0.3em}
        \textbf{Validation:}
        \begin{itemize}
            \item Per-iteration wall-clock time scales \textbf{linearly}
                with agent count across all three environments.
            \item Pearson $\rho > 0.999$ in every case.
        \end{itemize}

        \small This metric enables cross-configuration comparison in C1.
    \end{column}
    \begin{column}{0.46\textwidth}
        \begin{equation*}
            \text{Agent-steps} = |I| \times T
        \end{equation*}
        \begin{figure}
            \centering
            \includegraphics[width=\linewidth]{iter_cost.png}
            \caption{Mean wall-clock time per training iteration scales
                     linearly with agent count ($\rho > 0.999$).}
            \label{fig:agent-steps}
        \end{figure}
    \end{column}
\end{columns}

\note[note]{Smit only evaluated final performance}
\note[note]{not enough on costs and trade-offs on timing v performance}
\note[note]{
    Training steps is a common metric, 
    but does not consider cost of greater numbers of agents.
}
\note[note]{Wall-clock time would be a better proxy, but we use a cluster.}
\note[note]{Agent-steps gave us a fairer cost metric.}
\note[item]{
    This further allowed us to approximate the none-agent overhead cost in training.
}
\end{frame}


\begin{frame}{Waterworld}
    \begin{columns}
        \begin{column}{0.50\textwidth}
            \textbf{Environment character:}
            \begin{itemize}
                \item Continuous control, homogeneous structure.
                \item Low coordination demand; behavioral heterogeneity only.
            \end{itemize}
        \end{column}
        \begin{column}{0.45\textwidth}
            \begin{figure}
                \centering
                \includegraphics[width=0.5\linewidth]{waterworld.png}
                \caption{Waterworld environment~\cite{gupta2017}.}
                \label{fig:waterworld-env}
            \end{figure}
        \end{column}
    \end{columns}
\end{frame}

\begin{frame}{Waterworld Results}
Example Training Curves
\begin{columns}
    \begin{column}{0.48\textwidth}
        \begin{figure}
            \includegraphics[width=\linewidth]{Waterworld-4-agent.png}
            \caption{4-agent Waterworld: retrained vs. tabula rasa.}
        \end{figure}
    \end{column}
    \begin{column}{0.48\textwidth}
        \begin{figure}
            \includegraphics[width=\linewidth]{Waterworld-8-agent.png}
            \caption{8-agent Waterworld: retrained vs. tabula rasa.}
        \end{figure}
    \end{column}
\end{columns}

\note[note]{Note: Matching asymptotic performance}
\note[note]{But earlier convergent}
\note[note]{next we measure how much earlier}
\end{frame}


\begin{frame}{Waterworld Results}
\begin{columns}
    \begin{column}{0.48\textwidth}
        \textbf{Findings:}
        \begin{itemize}
            \item Pretraining \textbf{accelerates convergence}; benefit
                scales with the size ratio (pretrain $\rightarrow$ target).
            \item Final performance matches tabula rasa baselines.
        \end{itemize}

        \vspace{0.3em}
        \begin{alertblock}{Take-away}
            Clear curriculum win — low coordination means
            small-team policy transfers cleanly.
        \end{alertblock}
    \end{column}
    \begin{column}{0.50\textwidth}
        \begin{figure}
            \centering
            \includegraphics[width=\linewidth]{Waterworld-AUCs.png}
            \caption{Waterworld AUC improvement over tabula rasa.}
            \label{fig:waterworld-aucs}
        \end{figure}
    \end{column}
\end{columns}

\note[note]{Heatmap: compares area under the curve relative to tabula rasa}
\note[note]{rows are target team size}
\note[note]{columns, pre-training duration}
\note[note]{Clear gradient - larger teams get more benefit}
\note[note]{but overall, not too much pretraining needed}
\end{frame}

\begin{frame}{Multiwalker}
    \begin{columns}
        \begin{column}{0.50\textwidth}
            \textbf{Environment character:}
            \begin{itemize}
                \item Tight temporal coordination among physically coupled walkers.
                \item Fixed spatial roles create \textbf{static intrinsic heterogeneity}.
                \item Sparse rewards; single-agent failure terminates the episode.
            \end{itemize}
        \end{column}
        \begin{column}{0.45\textwidth}
            \begin{figure}
                \centering
                \includegraphics[width=0.75\linewidth]{multiwalker.png}
                \caption{Multiwalker environment~\cite{gupta2017}.}
                \label{fig:multiwalker-env}
            \end{figure}
        \end{column}
    \end{columns}
\end{frame}

\begin{frame}{Multiwalker Results}
Example Training Curves
\begin{columns}
    \begin{column}{0.48\textwidth}
        \begin{figure}
            \includegraphics[width=\linewidth]{Multiwalker-5-agent.png}
            \caption{5-agent Multiwalker: retrained vs. tabula rasa.}
        \end{figure}
    \end{column}
    \begin{column}{0.48\textwidth}
        \begin{figure}
            \includegraphics[width=\linewidth]{Multiwalker-6-agent.png}
            \caption{6-agent Multiwalker: retrained vs. tabula rasa.}
        \end{figure}
    \end{column}
\end{columns}

\note[note]{Multiwalker experienced drastically differing results.}
\note[note]{It became an issue with stability of training.}
\note[note]{Too many untrained agents caused failure too early for any to learn.}
\end{frame}

\begin{frame}{Multiwalker Results}
\begin{columns}
    \begin{column}{0.48\textwidth}
        \textbf{Findings:}
        \begin{itemize}
            \item For smaller teams (4--5), pretraining offers modest gains.
            \item For larger teams (6--7), tabula rasa \textbf{frequently fails}
                to converge — pretraining acts as \textbf{stabilization}.
            \item Pretraining provides baseline competence to survive early episodes.
        \end{itemize}

        \vspace{0.3em}
        \begin{alertblock}{Take-away}
            Task structure determines the \emph{nature} of the benefit:
            speed in Waterworld, stability in Multiwalker.
        \end{alertblock}
    \end{column}
    \begin{column}{0.50\textwidth}
        \begin{figure}
            \centering
            \includegraphics[width=\linewidth]{Multiwalker-AUCs.png}
            \caption{Multiwalker AUC improvement. Large teams
                     show dramatic gains where tabula rasa often fails.}
            \label{fig:multiwalker-aucs}
        \end{figure}
    \end{column}
\end{columns}

\note[note]{That stability chang is illsutrated here}
\note[note]{Small teams need less, training}
\note[note]{Larger teams benefit a lot}
\note[note]{Small teams show potential overfiting}
\end{frame}

\begin{frame}{Level-Based Foraging}
    \begin{columns}
        \begin{column}{0.50\textwidth}
            \textbf{Environment character:}
            \begin{itemize}
                \item Discrete gridworld; dynamic intrinsic heterogeneity.
                \item Higher level tasks require coordination.
                \item Shared observation space
            \end{itemize}
            \textbf{Size-invariant Observation Schemas Compared:}
            \begin{enumerate}
                \item Full: All teammate features included.
                \item Truncated: Fixed-size subset of teammate features.
                \item Ally-Ignorant: No teammate features included.
            \end{enumerate}
        \end{column}
        \begin{column}{0.45\textwidth}
            \begin{figure}
                \centering
                \includegraphics[width=0.75\linewidth]{lbf.png}
                \caption{Level-Based Foraging environment~\cite{papoudakis2021}.}
                \label{fig:lbf-env}
            \end{figure}
        \end{column}
    \end{columns}
    \note{
        The built-in observation space presented a problem that would inspire C2. \\
        Agents are passed each other's observations in a fixed concatenation.
    }
\note[note]{Explain the observation space schemas}
\end{frame}

\begin{frame}[shrink=20]{LBF: Schema Comparison Statistics}
    \vspace{2em}
\begin{columns}
    \begin{column}{0.65\linewidth}
        \begin{center}
            \scriptsize
            \begin{tabular}{cccccc}
                \toprule
                \textbf{Agents} & \textbf{Obs. 1} & \textbf{Obs. 2} & \textbf{t-stat} & \textbf{Bonf. p-value} \\
                \midrule
                3 & Ally-ignorant & Full      & 306.28 & \textbf{0.000} \\
                3 & Ally-ignorant & Truncated & 252.39 & \textbf{0.000} \\
                3 & Full          & Truncated &  -0.97 & 1.008 \\
                4 & Ally-ignorant & Full      &  53.09 & \textbf{0.000} \\
                4 & Ally-ignorant & Truncated &  88.01 & \textbf{0.000} \\
                4 & Full          & Truncated &   0.56 & 1.728 \\
                5 & Ally-ignorant & Full      &  73.35 & \textbf{0.000} \\
                5 & Ally-ignorant & Truncated &  23.91 & \textbf{0.000} \\
                5 & Full          & Truncated &  -1.49 & 0.423 \\
                6 & Ally-ignorant & Full      &  16.52 & \textbf{0.000} \\
                6 & Ally-ignorant & Truncated &  28.51 & \textbf{0.000} \\
                6 & Full          & Truncated &   1.47 & 0.440 \\
                7 & Ally-ignorant & Full      &  11.88 & \textbf{0.000} \\
                7 & Ally-ignorant & Truncated &  12.90 & \textbf{0.000} \\
                7 & Full          & Truncated &   0.45 & 1.955 \\
                \bottomrule
            \end{tabular}
        \end{center}
        \vspace{0.5em}
        \centering
        \footnotesize
        Pairwise $t$-tests comparing schema performance for each agent count. 
        Bonferroni correction applied within each group of comparisons by team size.
    \end{column}
    \begin{column}{0.30\linewidth}
        \begin{alertblock}{Take-away}
            Full and Truncated schema had a similar distribution of results,
            where Ally-ignorant was significantly different.
        \end{alertblock}
    \end{column}
\end{columns}

\note[note]{Bonferoni comparison, to determine statistical difference between schema}
\note[note]{Full and truncated knowledge may have been similar distributions}
\note[note]{but ally-ignorance was statistically different enough}
\end{frame}


\begin{frame}{Level-Based Foraging Results}
Example Training Curves
\begin{columns}
    \begin{column}{0.48\textwidth}
        \begin{figure}
            \includegraphics[width=\linewidth]{LBF-6-agent.png}
            \caption{6-agent LBF: retrained vs. tabula rasa.}
        \end{figure}
    \end{column}
    \begin{column}{0.48\textwidth}
        \begin{figure}
            \includegraphics[width=\linewidth]{LBF-7-agent.png}
            \caption{7-agent LBF: retrained vs. tabula rasa.}
        \end{figure}
    \end{column}
\end{columns}
\note{
    Modest early training advantages, very rapidly lost
}
\end{frame}

\begin{frame}[shrink=5]{Level-Based Foraging Results}
\vspace{2em}
\begin{columns}
    \begin{column}{0.48\textwidth}
        \textbf{Findings:}
        \begin{itemize}
            \item Benefits are \textbf{limited and inconsistent};
                gains diminish as target team grows.
            \item Observation schema matters critically.
            \item Dynamic heterogeneity creates coordination demands
                that \textbf{do not transfer} from smaller teams.
        \end{itemize}

        \vspace{0.3em}
        \begin{alertblock}{Take-away}
            When heterogeneity itself changes with team size,
            policy duplication cannot bridge the gap.
        \end{alertblock}
    \end{column}
    \begin{column}{0.50\textwidth}
        \begin{figure}
            \centering
            \includegraphics[width=\linewidth]{LBF-AUCs.png}
            \caption{LBF AUC improvement. Benefits favor small
                     target teams; diminish sharply at scale.}
            \label{fig:lbf-aucs}
        \end{figure}
    \end{column}
\end{columns}


\note[note]{small benefits, particular in larger teams}
\note[note]{likely an issue with coordination}
\note[note]{failure to specialize}
\end{frame}

\begin{frame}{C1 Summary \& Bridge to C2}
    \vspace{-0.3em}
    \begin{center}
    \renewcommand{\arraystretch}{1.5}
    \small
    \begin{tabular}{@{} l l l @{}}
        \toprule
        \textbf{Environment} & \textbf{Outcome} & \textbf{Driver} \\
        \midrule
        Waterworld      & Accelerated convergence
            & Low coordination; homogeneous structure \\
        Multiwalker     & Training stabilization
            & Sparse rewards; failure-sensitive coupling \\
        LBF             & Limited / inconsistent benefit
            & Dynamic heterogeneity; schema sensitivity \\
        \bottomrule
    \end{tabular}
    \end{center}

    \vspace{0.5em}
    \begin{alertblock}{Bridge to C2}
        If we design observation spaces to \emph{naturally support}
        structural heterogeneity, we eliminate these barriers at the source.
        \textbf{That is Contribution 2.}
    \end{alertblock}

\note{About 1.5 minutes.
    This is the synthesis slide — pull the three threads together.
    Walk through the table row by row (one sentence each).
    How the observation space changed between situations mattered}
\note[item]{In Waterworld it was not a problem (homogeneous structure).}
\note[item]{In Multiwalker it was manageable (fixed roles, compatible schemas).}
\note[item]{In LBF it was the primary obstacle (observation structure changes with team size).}
\end{frame}
